% =============================================================
%  appendice_stack.tex
%  Appendice A --- Lo stack delle chiamate
% =============================================================

\chapter{Lo stack delle chiamate}
\label{app:stack}

\begin{intuizione}
Capire come il computer gestisce le chiamate di funzione
è fondamentale per capire la tail recursion, lo stack overflow,
e perché certi programmi ricorsivi sono più efficienti di altri.
Questa appendice spiega il meccanismo dal basso.
\end{intuizione}

\vspace{4pt}

% ================================================================
\section{Cosa succede quando chiami una funzione}
% ================================================================

Quando un programma chiama una funzione, il computer deve:

\begin{enumerate}
  \item \textbf{Sospendere} l'esecuzione corrente
  \item \textbf{Ricordare} dove riprendere quando la funzione finisce
  \item \textbf{Passare} gli argomenti alla funzione
  \item \textbf{Eseguire} il corpo della funzione
  \item \textbf{Restituire} il risultato al chiamante
  \item \textbf{Riprendere} l'esecuzione dal punto sospeso
\end{enumerate}

\detail{Per fare tutto questo serve una struttura dati che tenga traccia
di tutte le chiamate attive nello stesso momento. Questa struttura
si chiama \textbf{call stack} (pila delle chiamate).}

\argomento{Il frame di attivazione}

\begin{concetto}{Stack frame (frame di attivazione)}
\keyline{}{Ogni chiamata di funzione crea un \emph{frame} sullo stack.}

\detail{Un frame contiene:}
\begin{itemize}
  \item I \textbf{parametri} passati alla funzione
  \item Le \textbf{variabili locali} definite nel corpo
  \item L'\textbf{indirizzo di ritorno}: dove riprendere dopo la chiamata
  \item Il \textbf{valore di ritorno} (o lo spazio per esso)
\end{itemize}

\detail{Quando la funzione termina, il suo frame viene \emph{rimosso}
(``fatto saltare'' --- \emph{popped}) dallo stack, e l'esecuzione
riprende dall'indirizzo di ritorno salvato nel frame.}
\end{concetto}

\begin{esempio}[Traccia dello stack per una chiamata semplice]
\begin{lstlisting}
(defun quadrato (x) (* x x))
(defun somma-quadrati (a b) (+ (quadrato a) (quadrato b)))

;; Chiamata: (somma-quadrati 3 4)
\end{lstlisting}

Lo stack durante l'esecuzione (cresce verso il basso):

\begin{tcolorbox}[
  enhanced, colback=black!3, colframe=black!25,
  arc=2pt, boxrule=0.4pt,
  left=10pt, right=10pt, top=8pt, bottom=8pt,
  before skip=6pt, after skip=6pt
]
\ttfamily\small
\begin{tabular}{|l|}
\hline
\textbf{frame: somma-quadrati} \\
\quad a = 3, b = 4 \\
\quad ritorna a: main \\
\hline
\textbf{frame: quadrato} \quad$\leftarrow$ stack top\\
\quad x = 3 \\
\quad ritorna a: somma-quadrati \\
\hline
\end{tabular}
\end{tcolorbox}

\detail{Quando \fn{quadrato} finisce, il suo frame viene rimosso.
Poi viene creato un nuovo frame per la seconda chiamata \fn{(quadrato 4)}.
Infine \fn{somma-quadrati} completa il calcolo e anche il suo frame
viene rimosso.}
\end{esempio}

% ================================================================
\section{La ricorsione e la crescita dello stack}
% ================================================================

\begin{concetto}{Ogni chiamata ricorsiva aggiunge un frame}
\keyline{}{Una funzione ricorsiva profonda N livelli\\
occupa N frame sullo stack \emph{contemporaneamente}.}

\detail{Questo è il punto critico: i frame non vengono rimossi finché
la ricorsione non comincia a ``tornare indietro''.
Tutti i frame intermedi devono rimanere in memoria perché
contengono il lavoro da fare ancora (la moltiplicazione, la somma, ecc.)
dopo che la chiamata ricorsiva restituisce il suo risultato.}
\end{concetto}

\begin{esempio}[Stack durante fattoriale ricorsivo]
\begin{lstlisting}
(defun fattoriale (n)
  (if (= n 0)
      1
      (* n (fattoriale (- n 1)))))  ; <-- il * aspetta il risultato

;; (fattoriale 4)
\end{lstlisting}

Stack al momento più profondo (quando si chiama \fn{(fattoriale 0)}):

\begin{tcolorbox}[
  enhanced, colback=black!3, colframe=black!25,
  arc=2pt, boxrule=0.4pt,
  left=10pt, right=10pt, top=8pt, bottom=8pt,
  before skip=6pt, after skip=6pt
]
\ttfamily\small
\begin{tabular}{|l|}
\hline
\textbf{frame: fattoriale}, n=4 \quad(deve ancora fare \texttt{* 4 ?}) \\
\hline
\textbf{frame: fattoriale}, n=3 \quad(deve ancora fare \texttt{* 3 ?}) \\
\hline
\textbf{frame: fattoriale}, n=2 \quad(deve ancora fare \texttt{* 2 ?}) \\
\hline
\textbf{frame: fattoriale}, n=1 \quad(deve ancora fare \texttt{* 1 ?}) \\
\hline
\textbf{frame: fattoriale}, n=0 \quad$\leftarrow$ stack top: ritorna 1 \\
\hline
\end{tabular}
\end{tcolorbox}

\detail{Tutti e cinque i frame esistono \emph{contemporaneamente} in
memoria. Ogni frame aspetta che quello sopra finisca per poter
completare la propria moltiplicazione.}
\end{esempio}

\argomento{Stack overflow}

\begin{concetto}{Stack overflow: quando lo stack esaurisce la memoria}
\keyline{}{Lo stack ha una dimensione massima fissa.}
\keyline{}{Se la ricorsione è troppo profonda, lo stack si esaurisce.}

\detail{Il sistema operativo alloca allo stack uno spazio limitato
(tipicamente 1--8 MB). Se il numero di frame supera quello che ci sta,
il programma va in crash con un errore chiamato \textbf{stack overflow}.}
\end{concetto}

\begin{esempio}[Stack overflow in pratica]
\begin{lstlisting}
;; Questa funzione va in stack overflow per n grande
(defun conta-giu (n)
  (if (= n 0)
      'fatto
      (conta-giu (- n 1))))   ; crea un frame per ogni chiamata

;; (conta-giu 100000)  --> potrebbe causare stack overflow
;; dipende dall'implementazione e dalla dimensione dello stack
\end{lstlisting}
\end{esempio}

\begin{confronto}
\textbf{Linguaggi diversi, limiti diversi.}
Java ha uno stack di default di 512KB--1MB: va in \texttt{StackOverflowError}
con poche migliaia di chiamate ricorsive.
Python ha un limite di ~1000 frame per default (configurabile con
\texttt{sys.setrecursionlimit}).
C/C++ dipende dall'OS: tipicamente 1--8MB di stack.
Common Lisp non ha un limite fisso per lo stack nei compilatori moderni
(come Allegro CL), ma la memoria disponibile è comunque finita.
\end{confronto}

% ================================================================
\section{Tail Call Optimization (TCO)}
% ================================================================

\begin{concetto}{L'intuizione della TCO}
\keyline{}{Se la chiamata ricorsiva è l'ultima cosa che la funzione fa,\\
il suo frame non serve più dopo la chiamata.}

\detail{Quando una funzione chiama sé stessa in \emph{posizione di coda}
(cioè: il risultato della chiamata ricorsiva è direttamente il risultato
della funzione, senza altri calcoli da fare), il frame corrente
può essere \emph{riutilizzato} invece di crearne uno nuovo.
Il compilatore lo trasforma essenzialmente in un \emph{salto} (jump)
invece di una \emph{chiamata} (call).}
\end{concetto}

\begin{esempio}[Non-tail vs tail: cosa cambia nello stack]
\begin{lstlisting}
;; NON tail-recursive: dopo la chiamata c'e' ancora da fare (* n ...)
(defun fattoriale-naif (n)
  (if (= n 0) 1
      (* n (fattoriale-naif (- n 1)))))
;; Stack: cresce di N frame -> possibile overflow

;; TAIL-recursive: la chiamata e' l'ultima operazione
(defun fattoriale (n)
  (labels ((aux (n acc)
             (if (= n 0) acc
                 (aux (- n 1) (* n acc)))))  ; niente dopo
    (aux n 1)))
;; Stack: il compilatore riusa lo stesso frame -> O(1) stack
\end{lstlisting}

Stack durante la versione tail-recursive (con TCO):

\begin{tcolorbox}[
  enhanced, colback=black!3, colframe=black!25,
  arc=2pt, boxrule=0.4pt,
  left=10pt, right=10pt, top=8pt, bottom=8pt,
  before skip=6pt, after skip=6pt
]
\ttfamily\small
\begin{tabular}{|l|}
\hline
\textbf{frame: aux} \quad(unico frame, aggiornato ad ogni passo) \\
\quad n=4, acc=1   $\to$ aggiorna a n=3, acc=4 \\
\quad n=3, acc=4   $\to$ aggiorna a n=2, acc=12 \\
\quad n=2, acc=12  $\to$ aggiorna a n=1, acc=24 \\
\quad n=1, acc=24  $\to$ aggiorna a n=0, acc=24 \\
\quad n=0, acc=24  $\to$ ritorna 24 \\
\hline
\end{tabular}
\end{tcolorbox}

\detail{Un solo frame per tutta la ricorsione, indipendentemente da N.
È equivalente a un ciclo \fn{while} come questo:}

\begin{lstlisting}
;; Equivalente imperativo:
;; n = N, acc = 1
;; while n > 0: acc = acc * n; n = n - 1
;; return acc
\end{lstlisting}
\end{esempio}

\begin{intuizione}
La tail recursion con TCO \textbf{è} un loop, visto da prospettive diverse.
Il compilatore trasforma la ricorsione in un salto all'inizio della
funzione con i parametri aggiornati --- esattamente come fa un ciclo.
Questo è il motivo per cui i linguaggi funzionali puri (Haskell, Erlang,
Scheme) possono fare a meno dei loop: la tail recursion con TCO è
altrettanto efficiente.
\end{intuizione}

% ================================================================
\section{TCO in Common Lisp}
% ================================================================

\begin{confronto}
\textbf{Lo standard ANSI Common Lisp non garantisce la TCO.}

\medskip
A differenza di Scheme (dove la TCO è obbligatoria per standard),
Common Lisp lascia la decisione all'implementazione.

\medskip
\textbf{In pratica:}
\begin{itemize}
  \item \textbf{Allegro CL, SBCL, CCL}: applicano TCO sulle chiamate
        in posizione di coda quando il codice è compilato
        (non necessariamente interpretato).
  \item \textbf{CLISP}: TCO limitata; per ricorsioni profonde
        preferire \fn{loop}/\fn{do}.
  \item \textbf{Regola pratica}: se la ricorsione può essere
        molto profonda (migliaia di livelli), scrivi in forma
        tail-recursive con \fn{labels} e compila il codice
        con \fn{compile} o caricandolo con \fn{:ld}/\fn{compile-file}.
\end{itemize}

\medskip
Nell'ambiente Allegro CL Free Express usato in questo corso,
la TCO è attiva sul codice compilato. Gli esempi con \fn{labels}
del Capitolo~3 sono sicuri.
\end{confronto}

\begin{nota}
Per verificare se la TCO è attiva in Allegro CL, puoi testare:
\begin{lstlisting}
;; Definisci e COMPILA (non solo valuta) la funzione
(defun conta (n)
  (labels ((iter (n)
             (if (= n 0) 'done (iter (- n 1)))))
    (iter n)))
(compile 'conta)          ; compila esplicitamente

;; Ora prova con un numero grande
(conta 10000000)          ; deve restituire DONE senza overflow
\end{lstlisting}
\end{nota}
